\section{Calibration--Ranking Separation}
\label{sec:calibration_ranking}

\subsection{Ordinal Classification with Missing Grades}

Let $\mathcal Y=\{0,\ldots,K-1\}$ be an ordered grade space.
Adaptation uses labeled images from the \emph{seen} grades $\mathcal O\subset\mathcal Y$, while the \emph{missing} grades $\mathcal U=\mathcal Y\setminus\mathcal O$ stay in the full-$K$ prediction space.
A score vector $Z(x)\in\mathbb R^K$ predicts $\hat y(x)=\arg\max_k Z_k(x)$.
For a missing grade $m$ and an adjacent seen grade $n$, the \emph{pair margin}
\begin{equation}
  g^{Z}_{mn}(x)=Z_m(x)-Z_n(x)
  \label{eq:pair_margin}
\end{equation}
ranks images drawn from the two grades.
Pair AUROC is the probability that an image of grade $m$ receives a larger margin than an image of grade $n$ (ties count one half), and MS-AUC averages it over adjacent missing--seen pairs.
Full-$K$ decisions and pairwise ranking are two properties of the same scores, and they can come apart.

\subsection{What Class-Bias Calibration Can and Cannot Change}

Post-hoc calibration adds an offset to the missing grades, $\operatorname{Cal}_c(Z)_k=Z_k+c\,\mathbf 1[k\in\mathcal U]$, as in calibrated stacking~\cite{chao2016gzsl}.
With a single missing grade $m$, the decision statistic is $Z_m(x)-\max_{k\neq m}Z_k(x)$, and varying $c$ moves its threshold without changing the ranking of images.
Calibration decides how many images qualify for the missing grade, not the order in which they qualify.
The pairwise ranking is equally immune to calibration, including temperature scaling~\cite{guo2017calibration}.

\begin{proposition}[Calibration preserves pairwise ranking]
\label{prop:cal}
Let $\widetilde Z(x)=aZ(x)+b$ with a shared scale $a>0$ and class biases $b\in\mathbb R^K$ independent of $x$.
For every pair $(m,n)$, $g^{\widetilde Z}_{mn}$ and $g^{Z}_{mn}$ induce the same image ranking, the same ties, and the same pair AUROC.
\end{proposition}

\noindent\emph{Proof.}
$g^{\widetilde Z}_{mn}(x)=a\,g^{Z}_{mn}(x)+b_m-b_n$ is strictly increasing in $g^{Z}_{mn}(x)$.\hfill$\square$

Fig.~\ref{fig:teaser}c shows the consequence on real data: recall rises, while the pair AUROC does not move.
To improve the order, the scores must change in an image-dependent way.

\subsection{Reference--Displacement Analysis}

We write the adapted scores as a zero-shot reference plus an adaptation displacement, $S(x)=T(x)+D(x)$ with $D=S-T$, so that every pair margin splits into two terms:
\begin{equation}
  g^{S}_{mn}(x)
  =\underbrace{T_m(x)-T_n(x)}_{\Delta T_{mn}(x)}
  +\underbrace{D_m(x)-D_n(x)}_{\Delta D_{mn}(x)}.
  \label{eq:margin_decomposition}
\end{equation}
The split is an identity, but it locates the problem.
The reference term $\Delta T_{mn}$ is fixed by the zero-shot model, so any ranking gain beyond zero-shot must come from $\Delta D_{mn}$.
For a missing grade, however, $D_m$ was never fitted to images of grade $m$.

A natural way to temper adaptation is to blend the adapted scores back toward the reference, $B_\lambda=T+\lambda D$, the output-space analogue of weight interpolation~\cite{wortsman2022robust}.
Blending rescales the displacement of every grade; with calibration on top, it still reaches only one family of rankings.

\begin{proposition}[Calibrated blending spans a one-parameter ranking family]
\label{prop:family}
Let $\widetilde B_\lambda=a(T+\lambda D)+b$ with $\lambda\ge0$, a shared scale $a>0$, and class biases $b\in\mathbb R^K$ independent of $x$.
For every pair $(m,n)$, $\widetilde B_\lambda$ ranks images exactly as $\Delta T_{mn}+\lambda\,\Delta D_{mn}$ does, so the ranking depends on $\lambda$ alone.
\end{proposition}

\noindent\emph{Proof.}
By \eqref{eq:margin_decomposition}, $g^{\widetilde B_\lambda}_{mn}=a(\Delta T_{mn}+\lambda\,\Delta D_{mn})+b_m-b_n$; apply Proposition~\ref{prop:cal}.\hfill$\square$

Blending thus traces a single ranking path through the zero-shot ($\lambda=0$) and adapted ($\lambda=1$) models, with $\lambda>1$ extrapolating.
Calibration leaves the ranking at each point of this path unchanged.
The strength $\lambda$ controls how much of the original displacement is kept, never where the missing grade's displacement comes from.
Section~\ref{sec:exp_matched} tests this family directly.
We go beyond this path by transferring displacement from the missing grade's seen neighbors (Sec.~\ref{sec:tie}).
